\chapter{Kernel Build System}
\label{ch:kernelbuild}

The Linux kernel has one of the most sophisticated build systems in open
source software. Based on GNU Make with a custom Kconfig language for
configuration, it supports cross-compilation, incremental builds, and
module management across 30+ architectures.

\section{Kbuild Architecture}

The kernel build system (Kbuild) is organized into several layers:

\begin{itemize}[nosep]
  \item \textbf{Top-level Makefile} (\texttt{Makefile}): Entry point.
        Sets architecture, cross-compilation flags, and delegates to the
        per-architecture Makefile.
  \item \textbf{Architecture Makefile} (\texttt{arch/x86/Makefile}):
        Architecture-specific build rules for x86\_64. Defines the
        bootable image target (\texttt{bzImage}).
  \item \textbf{Subdirectory Kbuild files} (\texttt{drivers/*/Kbuild}):
        Per-directory build rules. Each driver or subsystem has its own
        Kbuild file listing source files and their compilation flags.
  \item \textbf{Configuration system} (\texttt{Kconfig}): Menu-based
        configuration language. The kernel's \texttt{make menuconfig},
        \texttt{make defconfig}, and \texttt{make olddefconfig} targets
        all process Kconfig files.
\end{itemize}

\section{The Kconfig Language}

Kconfig is a domain-specific language for describing configuration options
and their dependencies. Each configuration option has:

\begin{lstlisting}
config NAMESPACES
    bool "Namespaces support"
    default y
    help
      Support for Linux namespaces, used for container isolation.

config CGROUPS
    bool "Control Groups support"
    select KERNFS
    help
      Support for process resource control groups.
\end{lstlisting}

Key Kconfig constructs:

\begin{longtable}{|p{3cm}|p{10cm}|}
\hline
\textbf{Keyword} & \textbf{Purpose} \\
\hline
\texttt{bool} & On/off option (y/n). \\
\texttt{tristate} & Three-state option (y/m/n). \\
\texttt{int} & Integer-valued option. \\
\texttt{string} & String-valued option (e.g., path to initramfs). \\
\texttt{default} & Default value if not otherwise specified. \\
\texttt{depends on} & Prerequisite option(s). \\
\texttt{select} & Force-enable another option. \\
\texttt{help} & Documentation text. \\
\texttt{prompt} & User-visible string in menuconfig. \\
\hline
\end{longtable}

\subsection{Dependency Resolution}

When option A depends on B, setting A to y forces B to y. The kernel's
\texttt{olddefconfig} target walks the entire dependency graph and
resolves all constraints:

\begin{lstlisting}[language=bash]
make olddefconfig
# Output: .config is updated with all dependencies satisfied
\end{lstlisting}

\section{Build Targets}

The most commonly used make targets:

\begin{longtable}{|p{3cm}|p{10cm}|}
\hline
\textbf{Target} & \textbf{Purpose} \\
\hline
\texttt{defconfig} & Generate default configuration for the target architecture. \\
\texttt{olddefconfig} & Update .config, resolving all dependency changes. \\
\texttt{menuconfig} & Interactive ncurses-based configuration editor. \\
\texttt{bzImage} & Build the compressed kernel image (x86\_64). \\
\texttt{modules} & Build all enabled loadable kernel modules. \\
\texttt{modules\_install} & Install modules to /lib/modules/\$(uname -r). \\
\texttt{clean} & Remove most build artifacts (keep .config). \\
\texttt{mrproper} & Remove all build artifacts including .config. \\
\texttt{cscope} / \texttt{tags} & Generate code navigation database. \\
\hline
\end{longtable}

\section{Incremental Builds}

The kernel build system tracks dependencies at the file level using
\%.o.cmd files. Each compiled object file has a corresponding .cmd file:

\begin{lstlisting}
# .built-in.a.cmd
cmd_built-in.a := rm -f built-in.a; \
    ar cDPrs built-in.a init/main.o arch/x86/init/init_task.o ...
\end{lstlisting}

When a source file changes, only the objects that depend on it are
recompiled. This makes development iteration fast:
\begin{lstlisting}[language=bash]
# After changing one driver source file:
make -j$(nproc) bzImage
# Only recompiles the changed file and relinks
\end{lstlisting}

\section{The Initramfs Integration}

When CONFIG\_INITRAMFS\_SOURCE is set, the build system:

\begin{enumerate}[nosep]
  \item Creates a temporary directory: \texttt{usr/initramfs\_data/}.
  \item Copies files listed in initramfs.list.
  \item Runs \texttt{scripts/gen\_initramfs\_list.sh} to generate
        \texttt{usr/initramfs\_data.cpio.gz}.
  \item Links the cpio archive into the kernel image via
        \texttt{init/initramfs.o}.
\end{enumerate}

The cpio archive is placed in the \texttt{.init.ramfs} ELF section:

\begin{lstlisting}[language=bash]
# Verify initramfs is embedded in kernel:
objcopy -O binary --only-section=.init.ramfs arch/x86/boot/bzImage /dev/stdout | cpio -t
\end{lstlisting}

\section{Cross-Compilation}

While Veilbox targets x86\_64, the kernel build system supports
cross-compilation for different architectures:

\begin{lstlisting}[language=bash]
# For ARM 32-bit:
make ARCH=arm CROSS_COMPILE=arm-linux-gnueabihf- defconfig
# For AArch64:
make ARCH=arm64 CROSS_COMPILE=aarch64-linux-gnu- defconfig
\end{lstlisting}

The architecture is selected by \texttt{ARCH} environment variable and
the toolchain prefix by \texttt{CROSS\_COMPILE}.

\section{Common Build Flags}

\begin{longtable}{|p{4cm}|p{4cm}|p{4cm}|}
\hline
\textbf{Flag} & \textbf{Typical Value} & \textbf{Purpose} \\
\hline
\texttt{-O2} & Kernel default & Optimization level. \\
\texttt{-fno-stack-protector} & (x86) & Disable stack canary. \\
\texttt{-mno-sse} & (x86) & Disable SSE instructions. \\
\texttt{-fno-PIE} & Kernel default & Disable position-independent. \\
\texttt{-march=x86-64} & (x86) & Target architecture. \\
\texttt{-mno-red-zone} & (x86\_64) & Disable red zone in ABI. \\
\hline
\end{longtable}

\section{Kbuild Variables Reference}

The kernel build system uses Make variables to control behavior:

\begin{longtable}{|p{4cm}|p{4cm}|p{4cm}|}
\hline
\textbf{Variable} & \textbf{Typical Value} & \textbf{Purpose} \\
\hline
ARCH & x86\_64 & Target architecture \\
CROSS\_COMPILE & (empty) & Toolchain prefix \\
KBUILD\_OUTPUT & (empty) & Separate build directory \\
KBUILD\_CFLAGS & (default) & Extra C compiler flags \\
KBUILD\_CPPFLAGS & (default) & Extra C preprocessor flags \\
KBUILD\_LDFLAGS & (default) & Extra linker flags \\
HOSTCC & gcc & Host compiler \\
HOSTCXX & g++ & Host C++ compiler \\
CC & gcc & Target compiler \\
LD & ld & Target linker \\
OBJCOPY & objcopy & Object copy tool \\
\hline
\end{longtable}

\section{Generating Patches for Submitting}

When modifying the Veilbox configuration or source, generate patches
for upstream submission:

\begin{lstlisting}[language=bash]
# Create a patch for kernel config changes
git diff kernel/configs/custom-os.config > veilbox-config.patch

# Create a patch for build script changes
git diff build.sh > build-script.patch

# Split patches for different topics
git format-patch -1 --stdout > single.patch
\end{lstlisting}

\section{Using External Modules}

The kernel supports building external modules outside the source tree:

\begin{lstlisting}[language=bash]
# Build an external module
make -C /path/to/linux-7.1.0-rc6 \
    M=/path/to/external/module \
    modules

# Install external module to initramfs
make -C /path/to/linux-7.1.0-rc6 \
    M=/path/to/external/module \
    modules_install \
    INSTALL_MOD_PATH=/path/to/rootfs
\end{lstlisting}
